% ============================================================
% Section 7: Discussion and Limitations
% ============================================================

\subsection{Implications for Resource-Constrained Deployment}
\label{sec:disc:deployment}

The primary practical implication of our results is that \emph{inference-time
orchestration architecture is a first-class engineering lever for improving
the capability of resource-constrained LLM deployments}. An organization
operating on a single 16\,GB GPU with a 7B model---a configuration costing
approximately \$1{,}000--\$2{,}000 in hardware---can achieve 2.21$\times$
higher code patch generation rates by adopting a hierarchical orchestration
structure rather than a monolithic single-agent loop, without any model
fine-tuning or additional hardware.

This has immediate implications for privacy-sensitive deployments (where
sending code to cloud APIs is prohibited), air-gapped environments (where
internet connectivity is absent), and cost-sensitive development workflows
(where per-token cloud API costs are prohibitive for high-volume agent use).

\subsection{Relationship to Scaling Laws}
\label{sec:disc:scaling}

Our result can be interpreted through the lens of the inference-time scaling
hypothesis: that additional compute at inference time can substitute for
additional parameters at training time~\cite{snell2024scaling,
brown2024large}. Hierarchical orchestration increases inference-time compute
in two ways: (1) by running multiple LLM calls per instance (one per agent
invocation), and (2) by generating significantly more tokens per instance
(7.7$\times$ more than the baseline). This additional compute is not
wasted---it is structural: it goes into targeted exploration (code localization),
verification (test re-execution), and plan synthesis (Domain Manager
summaries). Whether this compute investment is efficient relative to simply
using a larger model (e.g., a 70B model with equivalent VRAM via multi-GPU
setup) is an important open question.

\subsection{Limitations}
\label{sec:disc:limitations}

\paragraph{Metric.}
We report patch generation rate, not test-suite resolution rate. A system
that generates a syntactically valid non-empty patch that fails all tests
counts as a ``win'' under our metric. The gap between patch generation and
resolution is well-documented: \citet{yang2024sweagent} report that
SWE-agent generates patches on a substantially higher fraction of instances
than it resolves. Our resolution rate is thus likely lower than our $17.7\%$
patch generation figure. Future work should integrate the SWE-bench Docker
evaluation harness to report resolution rates.

\paragraph{Sympy exclusion.}
The 77 sympy instances constitute a systematic failure due to repository
setup, not model or orchestration capability. If sympy instances could be
properly evaluated, the headline numbers would shift. Our non-sympy rates
(H: 23.8\%, B: 10.8\%) may be a more representative characterization of
the system's capability on solvable instances.

\paragraph{Single model, single hardware configuration.}
All results use \texttt{Qwen2.5-Coder-7B-Instruct-AWQ} on an RTX 5080.
Generalizability to other model families (Code Llama, DeepSeek Coder,
StarCoder), quantization schemes (GPTQ, GGUF), and hardware configurations
(A100, H100, multi-GPU) is not established.

\paragraph{No ablation over hierarchy depth.}
We compare three-level (H) vs.\ one-level (B) directly, without evaluating
a two-level intermediate (orchestrator + flat workers, no domain managers).
Establishing the marginal contribution of the L2 Domain Manager layer over
a flat two-level hierarchy would require an additional experimental arm.

\paragraph{Determinism and variance.}
Each instance is run exactly once in each condition. Due to LLM sampling
non-determinism (temperature $> 0$ for most agents), there is instance-level
variance in outcomes. We do not report confidence intervals over multiple
runs due to the computational cost of re-running the full 300-instance
evaluation.

\paragraph{Wall-clock overhead.}
The hierarchical system requires an average of 94.8\,s per instance versus
5.2\,s for the baseline---an 18.2$\times$ wall-time overhead. For
interactive use cases where latency matters (e.g., IDE code completion),
this overhead is prohibitive. The hierarchical system is appropriate for
asynchronous batch workflows (e.g., nightly CI/CD bug triage, offline code
review) where throughput matters more than latency.
